\section{How the results will be described} \label{app:wording}

The outcome labels are computed by \code{s110\_final.py} and \code{s110\_localise.py}. The sentences that describe them
are produced by \code{s110\_summary.py} from their output files, and only these sentences describe the outcomes.
Combined positions, variants, later reprocessings or products, and individual transit times and depths appear only in
tables, never in these sentences. ``$p$'' is always the empirical $p$ of Appendix~\ref{app:localisation}.
``Recovered'' never means ``confirmed planet'' (Section~\ref{sec:limits}). A combined statement about A and B, or about
the system, is not part of this test; it is exploratory and cannot be cited for or against either outcome. In a dry
run, ``Sector 110'' is replaced by ``the archival dry run (Sector NN)''.

In these tables X stands for the signal (A or B), and the other letters (n, R, k, N, U, r, u, C, s, P, x) are filled in
by the script: for example n is the number of counted transits, R and r are false-detection levels with their trial
counts k of N, U and u their bounds, s the sensitivity, P the empirical $p$, and N in the localisation part the number of
calibration injections. Each sentence about signal X is built from three parts: the outcome (Table~\ref{tab:sent1}), the robustness
qualifier for O1 and O3 (Table~\ref{tab:sent2}), and the localisation for O1 and O2 (Table~\ref{tab:sent3}).

\begin{longtable}{>{\raggedright\arraybackslash}p{0.08\textwidth}>{\raggedright\arraybackslash}p{0.86\textwidth}}
\caption{First part of the sentence, by outcome.\label{tab:sent1}}\\
\toprule
Outcome & First part of the sentence \\
\midrule
\endfirsthead
\toprule
Outcome & First part of the sentence \\
\midrule
\endhead
O1 & ``Signal X was recovered in Sector 110 at the pre-registered times and depth (n counted transits; measured
  false-detection level for this product about R \% overall (k of N trials; binomial 95 \% upper bound U \%, approximate
  because the trials share data) and r \% for exactly n counted transits (k' of N'; binomial bound u \%))'', with the
  numbers of Table~\ref{tab:fa} for the product of the run. If no trial had exactly $n$ transits: ``; no trials with
  exactly n counted transits''. For a TESS-SPOC primary: ``; measured on SPOC 2-min data, no trials exist for the
  TESS-SPOC product''. \\
O2 & ``A dip was found at the predicted times of signal X, but it is not the predicted signal (C), as pre-registered;
  this contradicts the archival model of signal X'', with C the failed conditions as written by the script: ``shallower
  than predicted'', ``deeper than predicted'', ``timing off''. \\
O3 & ``The predicted signal X is excluded by Sector 110 (pre-registered sensitivity s $>$= 0.95)'' \\
O4 & ``Sector 110 cannot decide on signal X (reason: R).'' with R one of the fixed reasons written by the script: fewer
  than 10 usable cadences; no registered QLP flux column; no counted transit; noise pool has fewer than 100 windows; d not
  above det and not below lo; sensitivity below 0.95; both of the last two. \\
\bottomrule
\end{longtable}

\begin{longtable}{>{\raggedright\arraybackslash}p{0.38\textwidth}>{\raggedright\arraybackslash}p{0.56\textwidth}}
\caption{Added part for the robustness qualifier (O1 and O3 only).\label{tab:sent2}}\\
\toprule
Case & Added part \\
\midrule
\endfirsthead
\toprule
Case & Added part \\
\midrule
\endhead
at least one variant qualifies, and none contradicts & ``; robust'' \\
no variant reaches $\Sens \geq 0.95$ & ``; robustness not testable'' \\
either of the two above while variant (d) is pending & the same, followed by `` (provisional until 2027-02-28)'' \\
O1, and a qualifying variant returned O3 & ``; a processing variant with adequate sensitivity excluded it (variant x);
  this is reported and does not change the outcome'' \\
O3, and a qualifying variant returned O1 or O2 & ``; a processing variant with adequate sensitivity detected a dip
  (variant x); this is reported and does not change the outcome'' \\
\bottomrule
\end{longtable}

\begin{longtable}{>{\raggedright\arraybackslash}p{0.22\textwidth}>{\raggedright\arraybackslash}p{0.72\textwidth}}
\caption{Added part for the localisation (O1 and O2 only).\label{tab:sent3}}\\
\toprule
Case & Added part \\
\midrule
\endfirsthead
\toprule
Case & Added part \\
\midrule
\endhead
A, L1 & ``; its pixel position is compatible with TIC 4206066 (empirical p = P from N on-target injections, not below
  0.05).'' \\
A, L2 & ``; its pixel position is in tension with TIC 4206066 (empirical p = P from N on-target injections, below 0.05
  and not below 0.0027).'' \\
A, L3 & ``; its pixel position is off target (empirical p = P from N on-target injections, below 0.0027): the signal
  most likely comes from another source, and a planet on TIC 4206066 is disfavoured.'' \\
A, not possible & ``; its position could not be measured (reason as written by the localisation script).'' \\
B, measured & ``; B's pixel offset is reported in the table without a label and is not interpreted in this test.'' \\
B, not possible & ``; B's localisation was not possible (reason as written by the localisation script).'' \\
A or B, fallback primary & as ``not possible'', with the reason ``fallback product: localisation not possible'' \\
\bottomrule
\end{longtable}

For O3 and O4 no localisation is run, and the sentence ends with the robustness part or the reason. P is printed with
four significant digits; the stated relation to the thresholds is computed from the unrounded value.

Two cases are worded by the code slightly differently from Tables~\ref{tab:sent1} and \ref{tab:sent3}; as stated in
Section~\ref{sec:protocol}, the code defines the sentence. The O4 reason ``both of the last two'' is written as ``d not
above det and not below lo; sensitivity below 0.95''. If the localisation distance is not finite, the label for A
is \code{NOT\_POSSIBLE}; the sentence for A then ends ``its position could not be measured (label NOT\_POSSIBLE).'',
and the sentence for B gives the clause ``B's pixel offset is reported in the table without a label and is not
interpreted in this test.''.

Appendix~\ref{app:public} adds three things: for the failure cases that leave no output file (Appendix~\ref{app:p5}) the
same sentence parts are written from the posted files with the fixed reasons listed there; a result declared invalid
and not repairable gives the reason ``pre-registered procedure not followed'' (Appendix~\ref{app:p1}); and every summary
is posted with the line of variant outcomes of Appendix~\ref{app:p7}.
